% !TEX root = ../main.tex
\chapter{Rys historyczny i aktualny stan wiedzy}
\label{chap:stan-wiedzy}


\section{Historia rekonstrukcji scen trójwymiarowych}
\label{sec:historia}

Rekonstrukcja scen trójwymiarowych jest jednym z~problemów współczesnej wizji
komputerowej. Problem polega na odtworzeniu głębi z~płaskich zdjęć.
Zazwyczaj głębia ta jest tracona w~wyniku projekcji perspektywicznej.
Ewolucja tej dziedziny prowadzi nas od matematycznej fotogrametrii, przez
statystyczne dopasowywanie cech, aż po rewolucję, jaką przyniosło nam głębokie
uczenie maszynowe.

Podstawy rekonstrukcji scen 3D sięgają XIX wieku i~są związane
z~fotogrametrią. Fotogrametria to nauka zajmująca się określaniem kształtów
i~położeń obiektów na podstawie dostępnych zdjęć. Przed erą komputerów proces
ten opierał się na ręcznej identyfikacji punktów wspólnych na zdjęciach.
Automatyzacja nastąpiła wraz ze sformułowaniem matematycznych podstaw geometrii
wielowidokowej. Metody opracowane w~tym czasie pozwoliły powiązać ze sobą
punkty na różnych obrazach przedstawiających tę samą scenę.

Wczesne komputery nie posiadały jednak stabilnych metod automatycznego
znajdowania punktów wspólnych na zdjęciach. Kluczowa okazała się praca
Davida Lowe'a, która zaprezentowała algorytm SIFT \eng{Scale-Invariant Feature
Transform}~\cite{lowe2004sift}. Algorytm ten pozwolił na identyfikację punktów
charakterystycznych na obrazach w~sposób odporny na zmiany związane ze skalą
i~rotacją, a~także na częściowe zmiany oświetlenia i~perspektywy.

Połączenie deskryptorów lokalnych z~algorytmami, które były odporne na szum,
sprawiło, że można było stworzyć pierwsze w~pełni automatyczne systemy potoku SfM
\eng{Structure from Motion}. Photo Tourism~\cite{snavely2006phototourism} było
platformą, która pokazała, że możliwe jest zrekonstruowanie rzadkiej chmury
punktów oraz pozycji kamer z~tysięcy nieuporządkowanych zdjęć pobranych
z~internetu. Problem jednak polegał na tym, że SfM generowało wyłącznie rzadkie
chmury punktów, które, chociaż opisywały geometrię sceny z~perspektywy kamer,
były zbyt ubogie dla zastosowań wizualizacyjnych (VR, geologia). Rozwiązaniem
miało być opracowanie algorytmów MVS \eng{Multi-View Stereo}, które potrafiły
zagęścić chmurę punktów. Następnym krokiem była zamiana punktów na siatki
wielokątów. Połączenie SfM i~MVS było branżowym standardem. Metody te miały
jednak jedną dużą wadę: nie radziły sobie z~powierzchniami bezteksturowymi,
błyszczącymi, przezroczystymi czy zawierającymi powtarzalne wzory.

Powiew świeżości nastąpił w~2020 roku, kiedy zaczęły pojawiać się prace nad
NeRF \eng{Neural Radiance Fields}~\cite{mildenhall2020nerf}. Zamiast traktować
przestrzeń jako punkty czy trójkąty, badacze zaproponowali potraktowanie sceny
jako ciągłego pola funkcji, które jest reprezentowane przez głęboką sieć
neuronową (MLP). Chociaż NeRF zachwycał fotorealizmem, był bardzo kosztowny
obliczeniowo. To doprowadziło do narodzin 3D Gaussian Splatting
(3DGS)~\cite{kerbl2023gaussian}. Metoda ta łączyła zalety reprezentacji
jawnych i~ciągłych pól luminancji.


\section{Wprowadzenie do metody 3DGS}
\label{sec:3dgs}

\subsection{Reprezentacja sceny}
\label{subsec:reprezentacja}

Zaprezentowana w~2023 roku w~pracy~\cite{kerbl2023gaussian} metoda
3D Gaussian Splatting rozwiązywała praktycznie wszystkie problemy technologii
NeRF. Metoda ta porzuciła koncepcję ciągłego reprezentowania sceny przez sieć
MLP na rzecz reprezentacji za pomocą trójwymiarowych elipsoid Gaussa (splatów).
Scena składa się z~punktów, które stają się elipsoidami opisanymi przez zestaw
parametrów:
\begin{enumerate}
  \item Pozycja (środek): wektor $X = (x, y, z)$ określający punkt będący
        środkiem Gaussiana.
  \item Skala (wymiary): wektor $S = (s_x, s_y, s_z)$.
  \item Rotacja: reprezentowana przez kwaternion $Q = (q_w, q_x, q_y, q_z)$
        w~celu uniknięcia blokady przegubu.
  \item Nieprzezroczystość: wartość $\alpha \in [0, 1]$.
  \item Kolor zależny od kierunku: zestaw współczynników harmonik sferycznych,
        który pozwala na zmianę koloru elipsoidy w~zależności od kierunku.
\end{enumerate}

Kształt i~orientację Gaussiana opisuje macierz kowariancji
$\Sigma \in \R^{3 \times 3}$. Macierz ta musi być dodatnio określona, aby
zapewnić geometryczny sens, dlatego stosuje się następującą postać:
\begin{equation}
  \Sigma = R\,S\,S^{\top} R^{\top}.
  \label{eq:kowariancja}
\end{equation}
Macierz $R \in \mathrm{SO}(3)$ jest macierzą rotacji wyznaczaną z~kwaternionu
jednostkowego $Q$, a~$S$ jest macierzą diagonalną, która na przekątnej zawiera
wartości skali.

Popularnym formatem plików do przechowywania informacji o~Gaussianach jest
format PLY. Wartość skali jest w~nim przechowywana po zlogarytmowaniu, co
sprawia, że podczas odczytu z~pliku należy zastosować funkcję wykładniczą.
Nieprzezroczystość jest w~tym formacie przechowywana jako logit, zgodnie ze
wzorem:
\begin{equation}
  \alpha = \sigma(x) = \frac{1}{1 + \exp(-x)}.
  \label{eq:sigmoid}
\end{equation}
Kolor natomiast jest modelowany za pomocą harmonik sferycznych (SH) stopnia~3.
Wymaga to 16 współczynników SH dla każdego kanału RGB, co daje łącznie
48~wartości.

\subsection{Projekcja Gaussianów na ekran}
\label{subsec:projekcja}

Renderowanie sceny odbywa się poprzez rzutowanie trójwymiarowych Gaussianów
na ekran i~ich kompozytowanie. Podstawy tej operacji wywodzą się
z~pracy~\cite{zwicker2002ewa}, która zaprezentowała metodę EWA Splatting
\eng{Elliptical Weighted Average Splatting}.

Oznaczmy przez $W$ macierz transformacji widoku, która przekształca punkt
z~przestrzeni sceny do przestrzeni kamery. Projekcja perspektywiczna jest
przekształceniem nieliniowym, co uniemożliwia zastosowanie jej do Gaussianów.
Wykorzystuje się zatem jakobian $J$ do odwzorowania projekcji w~punkcie, który
jest reprezentowany przez wektor pozycji $X$. Macierz kowariancji Gaussiana
w~przestrzeni 2D przyjmuje postać:
\begin{equation}
  \Sigma' = J\,W\,\Sigma\,W^{\top} J^{\top}.
  \label{eq:kowariancja2d}
\end{equation}

Do obliczeń wykorzystuje się tylko górną lewą podmacierz macierzy $\Sigma'$.
Wtedy potrzebną nam macierz $2 \times 2$ możemy przedstawić jako
\begin{equation}
  \Sigma'_2 =
  \begin{bmatrix}
    a & b \\
    b & c
  \end{bmatrix},
  \label{eq:kowariancja2x2}
\end{equation}
gdzie $a$, $b$, $c$ to elementy macierzy kowariancji w~2D. Macierz ta
definiuje elipsę, której osie wyznaczone są przez wektory i~wartości własne
$\lambda_1$ i~$\lambda_2$ (służą do wyznaczania \textit{bounding box} dla
rasteryzatora). Przekształcenie to odpowiada rzutowi trójwymiarowego Gaussiana
na płaszczyznę ekranu z~uwzględnieniem perspektywy kamery.

\subsection{Renderowanie i sortowanie}
\label{subsec:renderowanie}

Finalna barwa piksela wyznaczana jest metodą \textit{alpha compositing}.
Oznacza to przetwarzanie Gaussianów w~kolejności od najdalszego do
najbliższego kamery. Kolor każdego piksela jest obliczany zgodnie
z~równaniem:
\begin{equation}
  C(p) = \sum_{i} c_i\, \alpha_i(p) \prod_{j<i} \bigl(1 - \alpha_j(p)\bigr).
  \label{eq:compositing}
\end{equation}
% TODO: iloczyn po j<i we wzorze (eq:compositing) odpowiada kolejności
% front-to-back (i = 1 to Gaussian NAJBLIŻSZY kamery). Tekst wyżej i niżej
% mówi o kolejności od najdalszego i numeracji malejąco wg głębokości —
% ujednolicić opis ze wzorem.
Indeks $i$ numeruje Gaussiany malejąco według głębokości \eng{depth}.
$c_i$ jest kolorem $i$-tego Gaussiana wyznaczonym przez SH dla bieżącego
kierunku patrzenia. Nieprzezroczystość $\alpha_i$ jest liczona w~zależności od
odległości piksela od centrum elipsy:
\begin{equation}
  \alpha_i(p) = o_i \exp\!\left(-\tfrac{1}{2}\, d^{\top} \left(\Sigma'_2\right)^{-1} d\right),
  \label{eq:alpha}
\end{equation}
% TODO: dopisać, czym jest o_i (nieprzezroczystość Gaussiana z p. 1.2.1)
% i mu'_2 (rzutowany środek Gaussiana).
gdzie $d = p - \mu'_2$ jest wektorem przesunięcia piksela od rzutowanego
centrum Gaussiana.

Sortowanie stanowi największe wyzwanie, jeśli chodzi o~wydajność metody.
Sortowanie na CPU może okazać się za mało efektywne, gdyż sceny potrafią
posiadać ponad 1~mln Gaussianów. Żeby rozwiązać ten problem, często stosuje się
algorytmy równoległe do sortowania na karcie graficznej. Najczęściej stosowane
jest sortowanie pozycyjne \eng{radix sort}. Implementacja sortowania
pozycyjnego na GPU pozwala zredukować czas sortowania miliona Gaussianów do
kilku milisekund, podczas gdy analogiczna operacja na CPU zajmuje czas rzędu
dziesiątek milisekund.

\subsection{Trening}
\label{subsec:trening}

Inicjalizacja zbioru Gaussianów zaczyna się od metody SfM, która generuje
rzadkie chmury punktów z~tych samych zdjęć użytych do treningu. Każdy punkt
chmury staje się środkiem jednego Gaussiana z~kulistą kowariancją i~losową
barwą SH.


Proces treningu minimalizuje funkcję straty, porównując obrazy wyrenderowane
ze zdjęciami ze zbioru treningowego. Funkcja straty to kombinacja błędu L1
i~miary strukturalnej D-SSIM.

Ważnym elementem treningu jest kontrola gęstości Gaussianów. Algorytm co pewną
liczbę iteracji dokonuje zagęszczania i~przycinania. Proces ten klonuje lub
dzieli Gaussiany o~dużym gradiencie pozycji, a~także usuwa Gaussiany o~małym
rozmiarze, niskiej nieprzezroczystości. Pozwala to na automatyczne
dostosowanie liczby i~rozmieszczenia Gaussianów w~zależności od sceny.

\section{Analiza istniejących rozwiązań}
\label{sec:analiza}

Istnieje wiele narzędzi, które umożliwiają wizualizację, edycję i~analizę
sceny za pomocą metody 3DGS. Rozwiązania te można podzielić na aplikacje
desktopowe wymagające CUDA, aplikacje webowe oraz implementacje oparte na
otwartych standardach i~bibliotekach graficznych. Poniżej omówiono trzy
reprezentatywne narzędzia.

\subsection{splatviz}
\label{subsec:splatviz}

Splatviz to narzędzie autorstwa Floriana Barthela. Jest to aplikacja
desktopowa napisana w~Pythonie, renderująca sceny 3DGS za pomocą kerneli CUDA
i~wyświetlająca interfejs za pomocą biblioteki ImGui osadzonej w~oknie OpenGL.
Każda funkcja analityczna to osobny widget UI.

Możliwości renderowania wykraczają poza standardowe narzędzia. Aplikacja ta
obsługuje tryb split-screen, co pozwala na porównanie wizualne między różnymi
metodami treningowymi lub parametrami. Są dostępne tryby debugowania, które
pokazują kanał głębokości i~alpha.

Tym, co wyróżnia splatviz, jest warstwa analityczna. Widget Evaluate pozwala
na wpisanie kodu w~Pythonie, który jest wykonywany po renderingu i~daje dostęp
do zmieniania kontekstu renderowania oraz wizualizacji zmian w~postaci
histogramów. Widget Edit pozwala na zmianę zbioru Gaussianów przez kod
w~Pythonie, co umożliwia filtrowanie według pewnych kryteriów.

Ograniczeniem tego rozwiązania jest konieczność kompilacji na urządzeniach
wspierających CUDA, co wyklucza użytkowników nieposiadających kart GPU NVIDIA.

\subsection{SuperSplat}
\label{subsec:supersplat}

SuperSplat firmy PlayCanvas to aplikacja webowa oparta na WebGL~2.0. Sprawia
to, że viewer ten działa w~przeglądarce bez instalacji. Jest to standard
branżowy w~zakresie prezentacji i~edycji scen 3DGS.

Rozwiązanie to posiada wiele funkcji edycyjnych: narzędzia do selekcji
i~usuwania Gaussianów, kompresję sceny i~kadrowanie obszaru roboczego oraz
eksport do formatów PLY lub SPLAT. Zawiera też wbudowaną historię zmian
z~funkcją cofania ostatnich zmian. Mamy też efekty post-processingu, wsparcie
WebXR, co pozwala na oglądanie scen na urządzeniach VR/AR, a~także do 25
adnotowanych hotspotów z~tytułami, opisami i~zapisanymi widokami kamery.
Funkcja Timeline pozwala na tworzenie animowanych przelotów przez scenę.

Tym, czego na pewno brakuje SuperSplat, są funkcje analityczne. Edytor
wyświetla histogramy podstawowych właściwości i~podstawowe statystyki
numeryczne sceny, ale nie posiada narzędzi analizy ani dostępu do parametrów
Gaussianów, które pozwalałyby na własne metryki jakości. Ograniczenia te
wynikają z~tego, że rozwiązanie to bazuje na WebGL~2.0, co utrudnia eksport
danych do warstwy analitycznej.

\subsection{3DGS.cpp}
\label{subsec:3dgscpp}

Projekt 3DGS.cpp autorstwa shg8 wykorzystuje Vulkan API i~compute pipeline.
Vulkan został wybrany ze względu na dostępność prymitywów warp-level zbliżonych
do CUDA. W~obecnej wersji API OpenGL jest uznawany za przestarzały na
platformach Apple. Obsługiwane platformy to Windows, Linux, macOS, iOS
i~visionOS.

Jest to aplikacja C++ z~modułową strukturą oddzielającą backend renderujący od
warstwy okienkowej. Shadery Vulkana (SPIR-V) implementują kolejne etapy
potoku: preprocessing Gaussianów, sortowanie pozycyjne na GPU i~rasteryzację.
Funkcje wizualizacyjne to m.in.\ kontrola kamery, ładowanie pliku PLY
i~podstawowe parametry renderowania.

Największa wartość rozwiązania polega na demonstracji tego, że można stworzyć
viewer, nie używając do tego CUDA, a~używając otwartych standardów. Projekt
udowadnia, że wydajność, którą zapewniają karty NVIDIA, jest osiągalna na
kartach AMD, Intel i~Apple Silicon, co rozszerza grupę odbiorców. Brakuje tutaj
natomiast funkcji analitycznych, edycyjnych i~trwałych operacji zapisu.

\subsection{Zestawienie i wnioski}
\label{subsec:zestawienie}

Tabela~\ref{tab:porownanie} zawiera zestawienie porównawcze omówionych
narzędzi według wybranych kryteriów funkcjonalnych i~technologicznych.

\begin{table}[htbp]
  \centering
  \caption{Zestawienie porównawcze narzędzi wizualizacji scen 3DGS}
  \label{tab:porownanie}
  \small
  \renewcommand{\arraystretch}{1.2}
  \begin{tabularx}{\textwidth}{
      >{\raggedright\arraybackslash}p{0.2\textwidth}
      >{\raggedright\arraybackslash}X
      >{\raggedright\arraybackslash}X
      >{\raggedright\arraybackslash}X}
    \toprule
    \textbf{Kryterium} & \textbf{splatviz} & \textbf{SuperSplat} & \textbf{3DGS.cpp} \\
    \midrule
    Stos technologiczny & Python + CUDA + ImGui & Web (PlayCanvas / WebGL~2) & C++ + Vulkan \\
    Platforma           & Linux/Win (NVIDIA) & Przeglądarka (uniwersalna) & Win/\allowbreak Linux/\allowbreak macOS/\allowbreak iOS \\
    Wizualizacja        & Zaawansowana (multi-scene, debug) & Dojrzała + WebXR & Bazowa \\
    Analityka           & Bardzo silna (Eval Widget, histogramy) & Podstawowa & Brak \\
    Edycja              & Minimalna (przez kod) & Bardzo silna (selekcje, timeline) & Brak \\
    Format wejściowy    & .ply, .yml & .ply, .splat, .ksplat, .sogs & .ply \\
    Zależność od CUDA   & Tak & Nie & Nie \\
    \bottomrule
  \end{tabularx}
\end{table}

Proponowany w~pracy viewer 3DGS będzie stworzony w~języku C++ z~użyciem
OpenGL. Będzie to rozwiązanie niewymagające CUDA, posiadające ergonomię
interfejsu oraz podstawowe narzędzia do analizy. W~potoku realizacji
zastosowane będą compute shadery (dostępne od OpenGL~4.3), żeby
zaimplementować sortowanie pozycyjne.
